%%%%%%%%%%  scan idx 52  |  folio 41 (continued)  %%%%%%%%%%

\underline{Teorema 1.11 (Schlessinger). Sea $\underline A$ un
grupoide cofibrado sobre $\underline C_\Lambda$. Entonces
$\underline A$ admite un objeto formal mínimamente versal si y solo si:}

%%%%%%%%%%  scan idx 53  |  folio 42  %%%%%%%%%%

(i) \underline{$\underline A$ es semihomogéneo;}

(ii) \underline{$[\underline A(k[\varepsilon])]$ es un espacio
vectorial de dimensión finita sobre $k$ (cf. 1.3(c)).}

\underline{Demostración. Necesidad.} Sea
$\xi\in\widehat{\underline A}$ un objeto formal mínimamente versal de
$\underline A$, y consideremos un diagrama
\[
\begin{tikzcd}[column sep=large,row sep=large]
a_1 \arrow[dr] & & a_2 \arrow[dl,"\alpha"] \\
& a &
\end{tikzcd}
\]
en $\underline A$, donde $\alpha$ es sobreyectivo. Como $\xi$ es
cuasiinicial, existe un morfismo $\xi\longrightarrow a_1$. Puesto que
$\alpha$ es sobreyectivo, el compuesto
$\xi\longrightarrow a_1\longrightarrow a$ se levanta a $a_2$; por
tanto, obtenemos un diagrama conmutativo
\[
\begin{tikzcd}[column sep=large,row sep=large]
& \xi \arrow[dl] \arrow[dr] & \\
a_1 \arrow[dr] & & a_2 \arrow[dl,"\alpha"] \\
& a &
\end{tikzcd}
\]
que demuestra la condición 1.2(1).

Demostremos ahora la condición 1.2(2). Sean
\[
\begin{tikzcd}[column sep=large]
& S \\
S\times_k k[\varepsilon]
 \arrow[ur,"\pi_1"] \arrow[dr,"\pi_2"'] & \\
& k[\varepsilon]
\end{tikzcd}
\]
las proyecciones, y sean
\[
\begin{tikzcd}[column sep=large]
a' \arrow[dr,"\alpha"'] & & b' \arrow[dl,"\beta"] \\
& a &
\end{tikzcd}
\]
$\pi_1$-morfismos en $\underline A$. Por la definición de $\xi$,
existe un diagrama conmutativo
\[
\begin{tikzcd}[column sep=large,row sep=large]
& \xi \arrow[dl,"u"'] \arrow[dr,"v"] & \\
a' \arrow[dr,"\alpha"'] & & b' \arrow[dl,"\beta"] \\
& a &
\end{tikzcd}
\]
y, por tanto, $\pi_1P(u)=\pi_1P(v)$. Supongamos ahora que
$(\pi_2)_*a'\simeq(\pi_2)_*b'$. La minimalidad de $\xi$ da
$\pi_2P(u)=\pi_2P(v)$, de donde

%%%%%%%%%%  scan idx 54  |  folio 43  %%%%%%%%%%

\[
 P(u)=\pi_1P(u)\times\pi_2P(u)
     =\pi_1P(v)\times\pi_2P(v)=P(v).
\]
Como $\underline A$ es un grupoide cofibrado, existe un único
isomorfismo $\sigma:a'\xrightarrow{\sim}b'$ en
$\underline A(S\times_k k[\varepsilon])$ tal que $\sigma u=v$.
La igualdad $\beta\sigma u=\alpha u$ implica entonces
$\beta\sigma=\alpha$.

\underline{Suficiencia.} Sean $c_1,\ldots,c_n$ objetos de
$\underline A(k[\varepsilon])$ cuyas clases forman una $k$-base de
$[\underline A(k[\varepsilon])]$. Elijamos
\[
 a_1\in\underline A(
 k[\varepsilon]\times_k\cdots\times_k k[\varepsilon])
\]
de modo que $(\pi_j)_*a_1\simeq c_j$. Entonces
\[
 \widetilde a_1(k[\varepsilon]):
 \operatorname{Hom}_{\Lambda\text{-alg}}(R_1,k[\varepsilon])
 \longrightarrow[\underline A(k[\varepsilon])]
\]
es biyectivo, donde $R_1=P(a_1)$. En particular,
$a_1\longrightarrow e$ es versal para las pequeñas prolongaciones de
$e$. Elijamos $a_2\longrightarrow a_1$, versal para las pequeñas
prolongaciones de $a_1$, de tal modo que
$\overline\Omega_{a_2}\longrightarrow\overline\Omega_{a_1}$ sea
una biyección. Continuemos por inducción, eligiendo
$a_{n+1}\longrightarrow a_n$ versal para las pequeñas prolongaciones
y que induzca una biyección sobre $\overline\Omega$, y pongamos
\[
 \xi=\varprojlim_n a_n.
\]
Entonces $\overline\Omega_\xi\longrightarrow\overline\Omega_{a_1}$
es biyectivo, por lo que
\[
 \xi(k[\varepsilon]):
 \operatorname{Hom}_{\Lambda\text{-alg}}(R,k[\varepsilon])
 \xrightarrow{\ \sim\ }
 [\underline A(k[\varepsilon])]
\]
es biyectivo.

Queda por demostrar la versalidad. Consideremos un diagrama
\[
\begin{tikzcd}
& a' \arrow[d,"\alpha"] \\
\xi \arrow[r,"h"] & a
\end{tikzcd}
\]
donde $\alpha$ es sobreyectivo. Debemos levantar $h$ a $a'$. Por
inducción, podemos suponer que $a'$ es una pequeña prolongación de
$a$. Como $P(a)$ es artiniano, $h$ se factoriza a través de algún
$a_q$,
\[
 \xi\longrightarrow a_q\longrightarrow a.
\]
La semihomogeneidad proporciona un diagrama conmutativo
\[
\begin{tikzcd}
& a'' \arrow[r] \arrow[d,"\alpha'"] & a' \arrow[d,"\alpha"] \\
\xi \arrow[r] & a_q \arrow[r] & a,
\end{tikzcd}
\]
donde $\alpha'$ es también una pequeña prolongación. Como
$a_{q+1}\longrightarrow a_q$ es versal para las pequeñas
prolongaciones de $a_q$, existe un diagrama conmutativo
\[
\begin{tikzcd}
& & a'' \arrow[d,"\alpha'"] \\
\xi \arrow[r] & a_{q+1} \arrow[r] \arrow[ur] & a_q,
\end{tikzcd}
\]
que proporciona el levantamiento requerido.

\underline{Observación 1.12.} Sea

%%%%%%%%%%  scan idx 55  |  folio 44  %%%%%%%%%%

\[
 \cdots\longrightarrow a_n\longrightarrow a_{n-1}
 \longrightarrow\cdots\longrightarrow a_2\longrightarrow a_1
 \longrightarrow a_0=e
\]
una sucesión infinita de flechas en $\underline A$ tal que
\[
\left\{
\begin{aligned}
&a_i\text{ es versal para las pequeñas prolongaciones de }a_{i-1},\\
&P(a_i)\longrightarrow P(a_{i-1})\text{ es sobreyectivo},\\
&\overline\Omega_{a_i}\longrightarrow
 \overline\Omega_{a_{i-1}}\text{ es biyectivo para }i\gg0.
\end{aligned}
\right.
\]
La demostración anterior muestra que $\varprojlim_i a_i$ es un objeto
formal versal de $\underline A$.

\underline{Observación 1.13.} Sea $\underline A$ un grupoide
cofibrado sobre $\underline C_\Lambda$. Si $\underline A$ satisface
1.11(i) y (ii), entonces también los satisface $[\underline A]$;
además, por 1.10(ii),
$\xi\in\operatorname{ob}(\widehat{\underline A})$ es un objeto formal
mínimamente versal de $\underline A$ si y solo si lo es de
$[\underline A]$.

\underline{Observación 1.14.} Sea $\underline A$ un grupoide
cofibrado sobre $\underline C_\Lambda$ que admite un objeto formal
mínimamente versal $\xi$. Todo invariante de $R=P(\xi)$, una
$\Lambda$-álgebra local completa de tipo formalmente finito sobre
$\Lambda$, es automáticamente un invariante de $\underline A$ sobre
$\underline C_\Lambda$. Así, por ejemplo, decimos que $\underline A$
es (formalmente) liso sobre $\underline C_\Lambda$ si $R$ es
(formalmente) liso sobre $\Lambda$, es decir, un álgebra de series
formales sobre $\Lambda$; y decimos que $\underline A$ es irreducible
sobre $\underline C_\Lambda$ si $\operatorname{Spec}(R)$ es
irreducible.

En la sección siguiente introducimos la lisura para un funtor
$\underline A\longrightarrow\underline B$ sobre
$\underline C_\Lambda$, donde $\underline A$ y $\underline B$ son
categorías cofibradas en grupoides sobre $\underline C_\Lambda$.
De 2.4 se deduce que el hecho de que $\underline A$ sea (formalmente)
liso sobre $\underline C_\Lambda$ equivale a que el funtor
$P:\underline A\longrightarrow\underline C_\Lambda$ sea liso, por
lo que no surge ambigüedad alguna. Definimos
\[
 \dim_{\underline C_\Lambda}\underline A
   =\dim(k\otimes_\Lambda R).
\]

Para ello, sea $\underline C_k$ la categoría de las $k$-álgebras
locales artinianas con cuerpo residual $k$. El funtor
$\underline C_k\longrightarrow\underline C_\Lambda$ es plenamente
fiel, y su adjunto izquierdo es
\[
 \underline C_\Lambda\longrightarrow\underline C_k,
 \qquad R\longmapsto k\otimes_\Lambda R.
\]
Sea $\underline A|_{\underline C_k}$ el producto fibrado de
$P:\underline A\longrightarrow\underline C_\Lambda$ por
$\underline C_k\longrightarrow\underline C_\Lambda$. Si $\xi$ es un
objeto formal mínimamente versal para $\underline A$ sobre
$\underline C_\Lambda$, entonces $k\otimes_\Lambda\xi$ es
mínimamente versal para $\underline A|_{\underline C_k}$ sobre
$\underline C_k$; la flecha
$\xi\longrightarrow k\otimes_\Lambda\xi$ es cocartesiana sobre
$R\longrightarrow k\otimes_\Lambda R$. En consecuencia,
\[
 \dim_{\underline C_\Lambda}\underline A
 =\dim_{\underline C_k}(\underline A|_{\underline C_k}).
\]

